% ============================================================
%  LIST OF SYMBOLS
%  Printed automatically in the Nomenclature, after the acronyms.
%  Acronyms live in docs/acronyms.tex; mathematical symbols live here.
%  Delete this file if the thesis does not need a list of symbols.
%
%  HOW TO USE
%  - Put every entry inside the symbollist environment, one per line:
%        $symbol$ & description \\
%  - Several related symbols can share one line, e.g. $T_1$, $T_2$.
%  - To group symbols by topic, start each group with
%        \symbolgroup{Group name}
%    and leave a little space before the next group by ending the
%    last entry of a group with \\[9pt] instead of \\
%  - Groups are optional: for a short list, just write the entries.
%  - Where a symbol follows the notation of a particular code or
%    reference (e.g. EN 1998-1), say so in the description.
%  - Text written before the list (here or in short subsections) can
%    explain general conventions, e.g. what a subscript denotes or
%    how to read a symbol that has more than one meaning.
% ============================================================

\begin{symbollist}

\symbolgroup{Ground motion and hazard}
$Sa(T)$ & spectral acceleration at period $T$ \\
$a_g(t)$ & ground motion time-history as a function of time $t$ \\
$H(IM > im)$ & hazard curve, the mean annual frequency of exceeding intensity $im$ \\[9pt]

\symbolgroup{Structural response and demand}
$\xi$ & damping ratio \\[9pt]

\symbolgroup{Fragility and risk}
$\eta$, $\beta$ & median and logarithmic dispersion of a fragility function \\

\end{symbollist}
